\clearpage
\section{Solutions to the exercises}
\label{sec-solutions-2-duality}

The common tool in these solutions is a change of viewpoint between a group
and its dual. We use it to count quadratic residues, organize Walsh and Haar
computations, and pass from finite Poisson summation to distributional and
sampling formulas.

\begin{correction}{exo-proof-geom-reciprocite}
\begin{enumerate}
\item For $1\leq b<p$, the representative of $(-1)^bb$ is $b$ if $b$ is
even and $p-b$ if $b$ is odd; it is therefore even in both cases.
If two representatives coming from $a,a'\in E$ coincide, then
$ra\equiv\pm ra'\pmod p$, so $a\equiv\pm a'\pmod p$.
The plus sign gives $a=a'$. For the minus sign, $a+a'$ would be $p$,
which is impossible because the sum is even. The map is injective, and
both sets have $(p-1)/2$ elements.
\item Taking products modulo $p$ gives
\[
 \prod_{b\in B_r}b\equiv r^{(p-1)/2}\prod_{a\in E}a,
 \qquad
 \prod_{a\in E}a\equiv(-1)^{\sum_{b\in B_r}b}\prod_{b\in B_r}b.
\]
All factors are nonzero modulo $p$, so they can be canceled.
Euler's criterion gives the assertion, first as a congruence and then as an
equality of signs, since $p$ is odd.
\item Sum $ra=p\lfloor ra/p\rfloor+[ra]_p$ over $E$.
The left-hand side is even, while $p$ is odd. Thus the sum of the remainders
and the sum of the quotients have the same parity.
\item A lattice point on the open segment would satisfy $py=rx$ with
$0<x<p$. Coprimality forces $p\mid x$, a contradiction.
For each integer $a\in\{1,\ldots,p-1\}$, the height $ra/p$ is not an integer,
so the positive integer ordinates below it are $1,\ldots,t_a$.
\item Since $ra/p$ is not an integer,
\[
 t_{p-a}=\left\lfloor r-\frac{ra}{p}\right\rfloor=r-1-t_a.
\]
The two terms have the same parity because $r-1$ is even.
Choose one representative $a\leq(p-1)/2$ from each pair $\{a,p-a\}$.
Exactly one member of the pair is even, so summing over the even members
or over these smaller representatives gives the same parity.
For such $a$, $ra/p<r/2$, hence $\sum_{a=1}^{(p-1)/2}t_a$ counts exactly
the interior lattice points of $AKH$.
This pairing is the needed counting argument; reflection across the oblique
diagonal is not a lattice-preserving symmetry in general.
\item Interchanging the coordinate axes identifies the analogous count for
$\legsymb pr$ with the interior lattice points of $AHL$.
There is no lattice point on their common diagonal, nor on the half-integer
outer edges $x=p/2$ and $y=r/2$. Together the two triangles contain
$(p-1)(r-1)/4$ lattice points. Consequently,
\[
 \legsymb rp\legsymb pr=(-1)^{(p-1)(r-1)/4}.
\]
\end{enumerate}
\end{correction}

\begin{correction}{exo-thm-fermat-finite-body}
\begin{enumerate}
\item The cyclic group $\FF_q^*=\langle g\rangle$ has a unique subgroup
of each order dividing $q-1$. If $k\mid q-1$, its $k$th powers are
$H_k=\langle g^k\rangle$, of order $(q-1)/k$.
\item Orthogonality in the quotient gives
$\sum_{i=0}^{k-1}\chi_i(x)=k\mathbf1_{H_k}(x)$ for $x\neq0$.
Multiplying by $\psi(x)$ and summing proves the formula.
If $\psi$ is nontrivial, the trivial multiplicative character contributes a
Gauss sum of modulus one, while the other $k-1$ sums have modulus $\sqrt q$.
Thus
\[
 |\wh{f_{H_k}}(\psi)|\leq\frac{1+(k-1)\sqrt q}{k}<\sqrt q.
\]
\item Each element of $H_k$ has exactly $k$ nonzero $k$th roots, so
$N=kN'$. Count $x+y+v=0$ with $v\in-H_k$.
The Fourier coefficients of $f_{-H_k}$ are those of $f_{H_k}$ with the
characters inverted; hence $\Phi(-H_k)=\Phi(H_k)$.
Exercise~\ref{exo-eqn-grpe-abelien-finite} yields
\[
 \left|N'-\frac{|A_1||A_2||H_k|}{q}\right|
 \leq\Phi(H_k)\sqrt{|A_1||A_2|}<\sqrt{|A_1||A_2|q}.
\]
Multiply by $k$ and use $k|H_k|=q-1$.
\item The main term in \eqref{eqn-fermat-majoration} exceeds its error bound
provided $k^2l_1l_2q^3<(q-1)^4$.
Under the proposed hypothesis,
\[
 k^2l_1l_2q^3\leq(q-4)q^3<(q-1)^4,
\]
because the difference on the right is $6q^2-4q+1>0$.
Thus $N>0$. For general $k$, in the cyclic group of order $q-1$,
$\langle g^k\rangle=\langle g^d\rangle$ for $d=\gcd(k,q-1)$.
Apply the result with $d\leq k$. A nonzero $d$th-power right-hand side is
also a $k$th power, which gives the required solution for the original exponent.
\item With $A_1=A_2=H_d$, we have $l_1=l_2=d\leq k$.
Thus $q\geq k^4+4$ implies the preceding criterion. All three entries are
nonzero powers, so taking their $k$th roots gives a solution of
\eqref{eq-thm-fermat-finite-body} in $(\FF_q^*)^3$.
\end{enumerate}
\end{correction}

\begin{correction}{exo-transforme-walsh-statistic}
The character indexed by the bit mask of $S$ is precisely
$\prod_{s\in S}s$ under the encoding $+=0$, $-=1$.
In natural binary order,
\[
 W_8f=(8\mu,4\mu_a,4\mu_b,4\mu_{ab},4\mu_c,4\mu_{ac},4\mu_{bc},4\mu_{abc})
     =(586,-92,20,6,-6,40,0,-2).
\]
An FWT followed by the indicated rescalings computes all quantities:
\[
\begin{array}{c|rrrrrrrr}
 &\mu&\mu_a&\mu_b&\mu_{ab}&\mu_c&\mu_{ac}&\mu_{bc}&\mu_{abc}\\\hline
 \text{cm}&73.25&-23&5&1.5&-1.5&10&0&-0.5
\end{array}
\]
For example, high lighting is associated with a mean height $23$ cm below
that at low lighting, after averaging over the other factors.
The relatively large $ac$ contrast indicates that the lighting contrast
changes with the fertilizer level. These are descriptive effects, without
an uncertainty estimate from this unreplicated table.
\end{correction}

\begin{correction}{exo-wavelet-haar}
\begin{enumerate}
\item Each wavelet has squared norm $2^j2^{-j}=1$ and integral zero.
Wavelets at the same scale have disjoint supports. At different scales,
the coarser wavelet is constant on the support of the finer one, so their
inner product vanishes. The constant $\psi_0$ is orthogonal to them all.
The first $2^J$ functions belong to $F_J$ and form an orthonormal family
with $\dim F_J=2^J$ members; they are therefore a basis.
\item For a dyadic cell $I$ of length $2^{-J}$, orthogonal projection onto
constants gives $f_J|_I=|I|^{-1}\int_I f$. Define the modulus of continuity
$\omega_f(h)=\sup_{|x-y|\leq h}|f(x)-f(y)|$. Then
\[
 |f_J(x)-f(x)|\leq\frac1{|I|}\int_I|f(y)-f(x)|\,\d y
                  \leq\omega_f(2^{-J}).
\]
The same bound holds at the endpoint under our convention.
Adding a wavelet coefficient at level $J$ replaces the average on its
support by the two averages on the half intervals. Distinct wavelets at
that level have disjoint supports. Every intermediate partial sum is
therefore an averaging operator on a partition whose largest cell has
length $2^{-J}$, and satisfies the same bound.
As $\omega_f(2^{-J})\to0$, all partial sums converge uniformly.
Continuous functions are dense in $L^2([0,1])$, so the span of the Haar
family is dense there. Together with orthonormality, this proves completeness.
\item The normalized indicators $\varphi_{j,k}$ are an orthonormal basis
of $F_j$. Directly on the two half intervals,
\begin{equation}
\label{eqn-recurrence-wavelets-haar}
 \varphi_{j,k}=\frac{\varphi_{j+1,2k}+\varphi_{j+1,2k+1}}{\sqrt2},
 \qquad
 \psi_{j,k}=\frac{\varphi_{j+1,2k}-\varphi_{j+1,2k+1}}{\sqrt2}.
\end{equation}
The $2^j$ wavelets at level $j$ are orthogonal to $F_j$ and belong to
$F_{j+1}$. Dimension counting gives the asserted orthogonal decomposition.
\item If $f$ has value $v_k$ on cell $k$, then
$x^{(0)}[k]=2^{-J/2}v_k$.
Each pair is transformed by
\[
 \frac1{\sqrt2}\begin{pmatrix}1&1\\1&-1\end{pmatrix}.
\]
This matrix is orthogonal and has determinant $-1$: it is a reflection,
not a rotation. Applying it to disjoint pairs and rearranging the outputs
is still orthogonal. Also,
\[
 \sum_kx^{(i)}[k]=2^{-i/2}\sum_kx^{(0)}[k],
 \qquad x^{(J)}[0]=2^{-J/2}\sum_kx^{(0)}[k].
\]
Since the input has $2^J$ entries, the last expression is
$2^{J/2}\operatorname{mean}(x^{(0)})$.
\item Take inner products with $f$ in
\eqref{eqn-recurrence-wavelets-haar} and proceed by induction on $i$.
The initial statement holds by definition of $x^{(0)}$.
Each stage produces the coefficients of $\varphi_{J-i,k}$ and
$\psi_{J-i,k}$. The final scaling function is $\varphi_{0,0}=\psi_0$.
The total number of input entries processed is
$N+N/2+\cdots+2=2N-2$, so the cost is $O(N)$.
\item Orthogonality at each stage gives
\[
 \norm{x^{(0)}}_2^2
 =\sum_{i=1}^J\norm{d^{(i)}}_2^2+|x^{(J)}[0]|^2
 =\norm{\Gamma x^{(0)}}_2^2.
\]
Thus $\Gamma^{\mathsf T}\Gamma=I$. Its rows, or equivalently the vectors
$\Gamma^{\mathsf T}e_m$, form the basis whose inner products are the output
coefficients. They are the discrete Haar wavelets, followed by the normalized
constant vector, in the specified packing. Fine wavelets are localized on
short dyadic blocks; Walsh functions have full support. Haar analysis costs
$O(N)$ operations, compared with $O(N\log N)$ for the FWT.
\end{enumerate}
Program~\ref{listing-transfo-haar} implements precisely these signs and
normalizations. The accompanying Python module also provides the inverse.
\begin{listing}[H]
\PythonCode{listing-transfo-haar}
\caption{Orthonormal Haar analysis, with details stored from fine to coarse.}
\label{listing-transfo-haar}
\end{listing}
\end{correction}

\begin{correction}{exo-compression-walsh}
\begin{enumerate}
\item For an $N\times N$ array $f$, the transform is $W_NfW_N^{\mathsf T}$,
with entries
\[
 \wh f[i,j]=\sum_{s,t=0}^{N-1}f[s,t]\chi_i(s)\chi_j(t).
\]
Since $W_N^2=NI$, the inverse applies the same transform and divides by $N^2$.
Transforming rows and then columns costs $O(N^2\log N)$.
Program~\ref{listing-fftw2d} also handles rectangular arrays when each side
length is a power of two; the inverse divisor is the total number of pixels.
\begin{listing}[H]
\PythonCode{listing-fftw2d}
\caption{A separable two-dimensional Walsh transform and a sequency approximation.}
\label{listing-fftw2d}
\end{listing}
\item Write $n_i^{(k)}$ for the sequency of row $i$ of $W_{2^k}$.
The two rows above index $i$ at the next level are concatenations
$(\chi_i,\chi_i)$ and $(\chi_i,-\chi_i)$.
Every row starts with $1$; let $\epsilon_i^{(k)}$ be zero or one according
as its last entry is $1$ or $-1$. Counting internal and middle sign changes,
\begin{equation}
\label{eqn-rec-num-change-signs}
 n_i^{(k+1)}=2n_i^{(k)}+\epsilon_i^{(k)},\qquad
 n_{i+2^k}^{(k+1)}=2n_i^{(k)}+1-\epsilon_i^{(k)}.
\end{equation}
The pair supplies exactly $2n_i^{(k)}$ and $2n_i^{(k)}+1$.
Starting from $n_0^{(0)}=0$, induction proves that the sequencies are all
distinct and exhaust $\{0,\ldots,2^k-1\}$.
\item The endpoint parities satisfy
$\epsilon_i^{(k+1)}=\epsilon_i^{(k)}$ and
$\epsilon_{i+2^k}^{(k+1)}=1-\epsilon_i^{(k)}$.
Program~\ref{listing-nbrchgtsigne} computes the sequencies from these
recurrences. Sorting them gives the permutation to apply to the spectrum.
For $N=8$, their natural-order values are $(0,7,3,4,1,6,2,5)$.
Small sequency means few changes along the ordered sample axis; small
Hamming degree instead counts the factors in a Boolean character.
\begin{listing}[H]
\PythonCode{listing-nbrchgtsigne}
\caption{The number of sign changes in each naturally ordered Walsh row.}
\label{listing-nbrchgtsigne}
\end{listing}
\item The approximation is the orthogonal projection onto the first $M$
sequency-ordered functions. Its normalized squared error is
\[
 \norm{f-f_M}_2^2=\frac1{N^2}\sum_{m=M}^{N-1}|\Ww_k(f)[i_m]|^2.
\]
Only the retained coefficients need be stored because their positions are
fixed by the ordering. Walsh reconstruction is piecewise constant and can
introduce visible steps. Fourier reconstruction uses smooth trigonometric
functions but can oscillate near discontinuities. The relative error depends
on the signal; neither basis is uniformly better. Finite-bit compression
also quantizes the retained coefficients.
\item If $s_i$ and $s_j$ are the vertical and horizontal sequencies,
respectively, order coefficient $(i,j)$ by the pair $(s_i+s_j,s_i)$.
Retain the first $M$ coefficients and invert the transform.
The function \texttt{walsh\_approx2d} in Program~\ref{listing-fftw2d}
implements this rule, resolving ties by vertical sequency.
\end{enumerate}
\end{correction}

\begin{correction}{exo-matrices-hadamard}
\begin{enumerate}
\item This is the identity $W_{2^k}W_{2^k}^{\mathsf T}=2^kI$.
\item Sign changes and permutations preserve pairwise row orthogonality.
Normalize the first row by column sign changes, then the first column by
row sign changes. Group columns by the sign pairs in rows two and three.
The counts obey
\[
 i+j+k+l=N,\quad i+j-k-l=0,\quad
 i-j+k-l=0,\quad i-j-k+l=0.
\]
Solving gives $i=j=k=l=N/4$. For orders below three, the matrices $(1)$ and
$W_2$ supply the two exceptions.
\item Since $\eta(-1)=-1$, we have $Q_{ji}=-Q_{ij}$.
For $c\neq0$, substitute $z=(b+c)/b$ on $\FF_p^*$ to obtain
\[
 \sum_{b\in\FF_p}\eta(b)\eta(b+c)
 =\sum_{z\in\FF_p\setminus\{1\}}\eta(z)=-1.
\]
For $c=0$ the sum is $p-1$, since the nonzero values of $\eta$ have square
one. Therefore the diagonal entries of $QQ^{\mathsf T}$ are $p-1$ and its
off-diagonal entries are $-1$, giving $pI_p-J$.
Each row and column of $Q$ sums to $\sum_b\eta(b)=0$, so $QJ=JQ=0$.
\item The off-diagonal blocks in $HH^{\mathsf T}$ vanish because
$Q\mathbf1=Q^{\mathsf T}\mathbf1=0$.
The top-left entry is $p+1$, and the bottom-right block is
\[
 J+(Q-I_p)(Q^{\mathsf T}-I_p)
 =J+(pI_p-J)-Q-Q^{\mathsf T}+I_p=(p+1)I_p.
\]
The two Python programs below construct $H$ and check its Gram matrix
using exact integer arithmetic; $p=7$ gives an order-eight example.
\item By the determinant bound from Gram--Schmidt, the volume is at most
the product of the column norms. Each norm is at most $\sqrt N$, proving
\eqref{eq-borne-hadamard}. For a Hadamard matrix,
$|\det H|^2=\det(HH^{\mathsf T})=N^N$.
Finally, compactness gives a maximizer on the entrywise cube.
With all other entries fixed, the determinant is affine in one entry, so
its absolute value is convex and attains a maximum at an endpoint of that
entry's interval. Move entries successively to $\pm1$ without decreasing
$|\det A|$. The resulting vertex is also a maximizer.
\end{enumerate}
\begin{listing}[H]
\PythonCode{listing-matrixpaley}
\caption{The Paley construction for a prime $p\equiv3\pmod4$.}
\label{listing-matrixpaley}
\end{listing}
\begin{listing}[H]
\PythonCode{listing-testpaley}
\caption{Verifying the Paley construction at order $32$.}
\label{listing-testpaley}
\end{listing}
\end{correction}

\begin{correction}{exo-tensor-product}
\begin{enumerate}
\item Block multiplication gives
$(A\otimes B)(C\otimes D)=AC\otimes BD$ whenever the dimensions match,
and $(A\otimes B)^*=A^*\otimes B^*$.
Thus $A^{\otimes n}(A^{\otimes n})^*=(AA^*)^{\otimes n}=s^nI_{s^n}$.
\item With $A=\bigl(\begin{smallmatrix}1&1\\1&-1\end{smallmatrix}\bigr)$,
the Walsh recursion gives $W_{2^n}=A^{\otimes n}$.
\item Split the input into $s$ consecutive blocks. Recursively apply
$A^{\otimes(n-1)}$ to each, then mix the resulting blocks by $A$.
Program~\ref{tensor-product-listing} implements this recursion.
For $N=s^n$, its cost is $O(sN\log_sN)$ with dense base multiplication,
or $O(N\log_sN)$ for fixed $s$. The inverse is
\[
 x=s^{-n}(A^*)^{\otimes n}y.
\]
Use the same algorithm with $A^*$ and divide the result by $s^n$.
\begin{listing}[H]
\PythonCode{tensor-product-listing}
\caption{Applying a tensor power without forming its matrix.}
\label{tensor-product-listing}
\end{listing}
\item The normalized base matrix $A_\alpha/\sqrt2$ is orthogonal, hence so
is its tensor power. At $\alpha=\pi/4$,
$U_\alpha=2^{-n/2}W_{2^n}$. At $\alpha=0$ it is diagonal, multiplying
entry $j$ by $(-1)^{w(j)}$, where $w(j)$ is its binary Hamming weight.
At $\alpha=\pi/2$ the base interchanges $0$ and $1$ in each coordinate.
Its tensor power complements every bit: $j\mapsto2^n-1-j$, which reverses
the sample vector. It does not reverse the order of bits within $j$.
\end{enumerate}
\end{correction}

\begin{correction}{exo-generalization-mac-williams}
\begin{enumerate}
\item If $a\neq0$, choose a coordinate $i$ with $a_i\neq0$ and set
$x=a_i^{-1}e_i$. Then $\dotp ax=1$. The linear map
$a\mapsto\dotp a{\cdot}$ is injective between spaces of dimension $k$,
so it is an isomorphism onto $E^*$.
\item The character map is a group homomorphism. If $a\neq0$, the same
coordinate choice lets $\dotp ax$ take every value of $\FF_q$.
The trace is surjective onto $\FF_p$, so some $x$ satisfies
$\operatorname{Tr}(\dotp ax)=1$ and $\chi_a(x)=e^{2\pi\imath/p}\neq1$.
Thus the map is injective; both groups have $q^k$ elements, so it is bijective.
The character $\psi_1$ itself need not be injective when $q$ is not prime.
\item If $a\in H^\perp$, then clearly $\chi_a(h)=1$ on $H$.
Conversely, if $\chi_a(h)=1$ for every $h\in H$, scalar closure gives
\[
 \operatorname{Tr}(\lambda\dotp ah)=0
 \quad\text{for every }\lambda\in\FF_q\text{ and }h\in H.
\]
The trace pairing is nondegenerate, so $\dotp ah=0$ for every $h$.
Without scalar closure this inference fails. For example, the subfield
$\{0,1\}\subset\FF_4$ is an additive subgroup but not an $\FF_4$-linear
subspace: multiplication of $1$ by an element outside $\FF_2$ leaves it.
\item Apply finite Poisson summation to
$f(t)=x^{k-w(t)}y^{w(t)}$. The transform factors over coordinates.
At a coordinate with $a_i=0$, the factor is $x+(q-1)y$.
At a coordinate with $a_i\neq0$, character orthogonality gives
$\sum_{t_i\neq0}\psi_1(a_it_i)=-1$, so the factor is $x-y$.
Therefore
\[
 \wh f(\chi_a)=\bigl(x+(q-1)y\bigr)^{k-w(a)}(x-y)^{w(a)}.
\]
Since $|H^\perp|=q^k/|H|$, Poisson summation over $H^\perp$ gives the
claimed identity for every $x,y\in\CC$, and hence as a polynomial identity.
\end{enumerate}
\end{correction}

\begin{correction}{exo-formula-fish-distrib}
For a Schwartz function $\varphi$, the series
$\sum_n|\varphi(ns)|$ is bounded by a constant times the Schwartz seminorm
$\sup_x(1+|x|)^2|\varphi(x)|$. Thus $\Pi_s$ is tempered.
Under the stated Fourier convention, rescaled Poisson summation gives
\[
 \sum_{n\in\ZZ}\wh\varphi(ns)
       =\frac{2\pi}s\sum_{k\in\ZZ}\varphi(2\pi k/s).
\]
Both sums converge absolutely. By the definition of the distributional
transform, the left side is
$\langle\Pi_s,\wh\varphi\rangle=\langle\wh{\Pi_s},\varphi\rangle$;
the right side is $\langle(2\pi/s)\Pi_{2\pi/s},\varphi\rangle$.
Equality for all Schwartz test functions proves the formula.
\end{correction}

\begin{correction}{exo-sample-shannon}
\begin{enumerate}
\item Compact support and Cauchy--Schwarz imply $\wh f\in L^1$.
Fourier inversion therefore gives the continuous representative
\[
 f(t)=\frac1{2\pi}\int_{I_T}\wh f(\omega)e^{\imath\omega t}\,\d\omega.
\]
Multiplication by any power of $\omega$ preserves integrability on $I_T$,
so differentiation under the integral proves smoothness. Two continuous
representatives equal almost everywhere must agree everywhere.
Cauchy--Schwarz and Plancherel also give
\[
 |f(t)|\leq\frac1{2\pi}\sqrt{\frac{2\pi}T}\norm{\wh f}_2
         =T^{-1/2}\norm f_2.
\]
\item The functions $e_n(\omega)=\sqrt{T/(2\pi)}e^{-\imath nT\omega}$
form an orthonormal basis of $L^2(I_T)$.
Their coefficients are
$\dotp{\wh f}{e_n}=\sqrt{2\pi T}\,f(nT)$, so
\[
 \wh f(\omega)=T\sum_{n\in\ZZ}f(nT)e^{-\imath nT\omega}
 \quad\text{in }L^2(I_T),\qquad
 \norm f_2^2=T\sum_{n\in\ZZ}|f(nT)|^2.
\]
Extend this series periodically with period $2\pi/T$.
It defines the periodization of $\wh f$, since $\wh f$ vanishes outside
one period up to null sets. On the other hand, the transform of
$f_d=\sum_nf(nT)\delta_{nT}$ is the series
$\sum_nf(nT)e^{-\imath nT\omega}$ in tempered distributions.
The $\ell^2$ bound on the samples justifies this distributional series.
Comparing gives the stated periodization identity; in the interior of the
base interval only its unshifted term is nonzero almost everywhere.
This argument also supplies the convergence needed to use the sampling identity.
\item Direct integration gives
$\Ff^{-1}(\mathbf1_{I_T})=T^{-1}\sinc_T$.
Invert the $L^2(I_T)$ Fourier series from the previous question, extending
each term by zero outside $I_T$. This gives
\eqref{eqn-formula-reconstruction-shannon} in $L^2(\RR)$.
Every finite sum and the remainder belong to $E_T$; the bound from the
first question makes $L^2$ convergence imply uniform convergence.
\item The transform of $g_n=\sinc_T(\cdot-nT)$ is
$T\mathbf1_{I_T}(\omega)e^{-\imath nT\omega}$.
Plancherel gives $\dotp{g_n}{g_m}=T\delta_{nm}$, and completeness follows
from the exponential basis. Thus $g_n/\sqrt T$ is orthonormal and complete.
For arbitrary $g\in L^2(\RR)$, the orthogonal projection is
\[
 P_{E_T}g=\Ff^{-1}(\mathbf1_{I_T}\wh g)
         =\frac1T\sum_{n\in\ZZ}\dotp g{g_n}\,g_n.
\]
When $f\in E_T$, inversion also gives $\dotp f{g_n}=T f(nT)$.
The factor $T$ is essential unless the sample spacing is one.
\end{enumerate}
\end{correction}
